\section{Method}
\label{sec:method}

\subsection{Evaluation protocol}
\label{sec:protocol}

Each grain is a single contiguous series.
We cut it into train, validation, and test blocks (Table~\ref{tab:protocol}).
The validation block is used only to choose hyperparameters and thresholds.
When an estimator can be refit, we refit it on train and validation together and then score the untouched test block.
The official rank is raw test sMAPE.
Lower is better.
Root mean squared error is a secondary descriptor of scale.
We do not average ranks across grains.
A Kernel Ridge fit with test sMAPE at least 180 is treated as degenerate and is removed from the rank.

\begin{table}[t]
  \caption{Evaluation protocol by grain. The seasonal reference is the denominator of the skill score.}
  \label{tab:protocol}
  \centering
  \small
  \begin{tabular}{llll}
    \toprule
    Grain & Blocks & Official rank & Seasonal reference \\
    \midrule
    Monthly & train $|$ val(12) $|$ test(12) & sMAPE (RMSE secondary) & lag 12 \\
    Daily A (calendar) & train $|$ val(90) $|$ test(90) & sMAPE (RMSE secondary) & lag 7 \\
    Daily B (business days) & train $|$ val(90) $|$ test(90) & sMAPE (RMSE secondary) & lag 5 \\
    Weekly & train $|$ val(13) $|$ test(13) & sMAPE (RMSE secondary) & lag 52 \\
    \bottomrule
  \end{tabular}
\end{table}

Let $\epsilon$ be a small positive tolerance.
MAPE is averaged on $\mathcal{T}_{\mathrm{MAPE}} = \{t : |y_t| \gt \epsilon\}$, because a zero collection day makes the percentage undefined.
sMAPE is averaged on $\mathcal{T}_{\mathrm{sMAPE}} = \{t : |y_t| + |\hat y_t| \gt \epsilon\}$, which drops only the null-denominator case in which both the outcome and the forecast are zero.
RMSE uses every test time.
With $n_{\mathrm{M}} = |\mathcal{T}_{\mathrm{MAPE}}|$ and $n_{\mathrm{S}} = |\mathcal{T}_{\mathrm{sMAPE}}|$,
\begin{align}
  \mathrm{MAPE}
  &= \frac{100}{n_{\mathrm{M}}}
     \sum_{t \in \mathcal{T}_{\mathrm{MAPE}}}
     \left|\frac{y_t - \hat y_t}{y_t}\right|,
  \label{eq:mape}\\
  \mathrm{sMAPE}
  &= \frac{100}{n_{\mathrm{S}}}
     \sum_{t \in \mathcal{T}_{\mathrm{sMAPE}}}
     \frac{2\,|y_t - \hat y_t|}{|y_t| + |\hat y_t|},
  \label{eq:smape}\\
  \mathrm{RMSE}
  &= \sqrt{\frac{1}{n}\sum_{t=1}^{n}(y_t - \hat y_t)^2}.
  \label{eq:rmse}
\end{align}
Mean absolute error and a weighted absolute percentage error are implemented beside these scores and are not used to rank models in this freeze.
On the monthly panel, exact zeros are rare.
On calendar-daily slips they dominate the sample, so sMAPE and RMSE are the scores that remain interpretable if MAPE has dropped most of the week.

\subsection{Skill against the seasonal naive}
\label{sec:ss}

For a candidate forecast with sMAPE equal to $A$ and a reference forecast with sMAPE equal to $A_{\mathrm{ref}}$,
\begin{equation}
  \mathrm{SS}_{\mathrm{sMAPE}}
  = 1 - \frac{A}{A_{\mathrm{ref}}},
  \label{eq:ss}
\end{equation}
where $A_{\mathrm{ref}}$ is the test sMAPE of the seasonal naive at the grain lag in Table~\ref{tab:protocol}.
A positive score means lower sMAPE than that seasonal naive.
The seasonal naive itself scores zero.
A negative score means higher sMAPE than the seasonal naive.
The score is left undefined when $A_{\mathrm{ref}}$ is near zero or is not finite.
The reference is always the seasonal naive of the grain under study.
On business days that naive is lag~5.
It is a weak reference on this series, because a lag-1 naive already beats it.
Positive skill on Ranking~B is therefore expected for many families and is not evidence of an error in~\eqref{eq:ss}.

\subsection{Stage 1: point-forecast families}
\label{sec:m0m9}

Stage~1 is a ladder from seasonal benchmarks to a zero-shot foundation model (Table~\ref{tab:ladder}).
The same specification is applied at every grain.
Only the seasonal lag, the covariate set, and a small number of length-dependent caps change.
Tabular regressors (ordinary least squares, ridge, histogram gradient boosting, and the adaptive forest) produce one-step walk-forward forecasts: at each test time the features use information available before that time, and the fitted map is applied to that row.
The perceptron is also scored one step ahead, feeding observed lags rather than its own previous forecast.

\begin{table}[t]
  \caption{Stage~1 and Stage~2 families. Selection, when a family has hyperparameters, minimizes validation sMAPE.}
  \label{tab:ladder}
  \centering
  \small
  \begin{tabular}{llp{6.4cm}}
    \toprule
    ID & Family & Role in this protocol \\
    \midrule
    M0 & Naive and seasonal naive & Copy $y_{t-k}$; skill-score reference \\
    M1 & OLS and ridge & Lag regression; ridge penalty on a validation grid \\
    M2 & SARIMAX & Short order grid by AIC; multi-step test forecast \\
    M3 & Histogram gradient boosting & Lag features; walk-forward \\
    M4 & Prophet & Trend plus seasonal Fourier terms \\
    M5 & Multilayer perceptron & Two hidden layers on standardized lags \\
    M6 & Adaptive random forest & Online forest, regression form \\
    M7 & Gaussian process & Kernel and noise chosen on validation \\
    M8 & Kernel ridge & RBF or Laplacian kernel; degenerate fits dropped \\
    M9 & TimesFM~3.0 & Zero-shot; no local gradient step \\
    E1 & Selective TEM & Energy threshold; abstain to seasonal naive \\
    E2 & Disagreement proxy & Switch between boosting and seasonal naive \\
    E3 & Residual diffusion & Reduced TimeGrad-style sampler \\
    E4 & Score sampler & Denoising score matching and Langevin \\
    \bottomrule
  \end{tabular}
\end{table}

\paragraph{M0: naive and seasonal naive.}
The forecast is the copy $\hat y_t = y_{t-k}$~\citep{hyndman2021fpp3}.
We report $k = 1$ and the seasonal lag of the grain.
Calendar-daily fits also report $k = 365$.
No parameter is estimated.
The seasonal member is the reference in~\eqref{eq:ss}: lag~12 (monthly), lag~7 (calendar daily), lag~5 (business daily), and lag~52 (weekly).

\paragraph{M1: ordinary least squares and ridge.}
The tabular forecast is the linear model $y_t = x_t^{\top}\beta + \varepsilon_t$.
On the monthly panel, $x_t$ stacks the collection and stock lags in Section~\ref{sec:data}.
On the other grains, $x_t$ stacks the target lags, slip-count lags, and calendar fields of that grain.
Ridge replaces the residual sum of squares by
\begin{equation}
  \hat\beta_{\alpha}
  = \arg\min_{\beta}
    \|y - X\beta\|_2^2 + \alpha\|\beta\|_2^2,
  \label{eq:ridge}
\end{equation}
with $\alpha \in \{0.1, 1, 10\}$ chosen by walk-forward validation sMAPE~\citep{hoerl1970ridge}.
Ordinary least squares is the unpenalized counterpart.
Both are refit in a one-step walk-forward pass on the test block.

\paragraph{M2: SARIMAX.}
We fit a seasonal ARIMA specification $(p,d,q)\times(P,D,Q)_s$ in the Box--Jenkins sense~\citep{box1970tsa,hyndman2021fpp3}.
The seasonal period is $s = 12$ on the monthly panel and $s = 7$ on calendar daily data.
On the weekly series the fitted seasonal period is $s = 4$, a short within-month cycle.
A period of 52 is not identified on a series whose test block is only 13 weeks.
The skill-score reference on that grain remains the lag-52 naive, not the SARIMAX seasonal period.
A short grid of orders is scored by AIC on the training block.
The selected specification is refit on train and validation and then forecasts the test horizon.
Prediction intervals from the Gaussian state-space forecast are available and are not part of the rank.

\paragraph{M3: histogram gradient boosting.}
Histogram gradient boosting is a histogram-split form of Friedman's gradient boosting machine, fit on the same lag features as M1~\citep{friedman2001greedy}.
The grid covers the iteration budget, the learning rate, and the tree depth.
Each candidate is scored by walk-forward validation sMAPE, with early stopping inside the booster.
The selected booster is walked forward through the test block.

\paragraph{M4: Prophet.}
Prophet decomposes the series into a piecewise-linear trend, additive Fourier seasonality, and a holiday term if one is supplied~\citep{taylor2018prophet}.
We grid the changepoint prior scale against the seasonality prior scale and keep the pair with the best validation sMAPE.
Monthly fits use yearly seasonality.
Daily fits use weekly and yearly seasonality.
No holiday calendar is passed in, including on Ranking~B.
The selected model is refit on train and validation and forecasts the test horizon in one shot.

\paragraph{M5: multilayer perceptron.}
M5 is a feed-forward network with two hidden layers and ReLU activations, mapping a window of standardized lags of $y$ to the next value~\citep{goodfellow2016dl}.
It is a lag-window perceptron.
It is not an N-BEATS basis expansion.
A standard scaler is fit on the training block only.
Hidden widths are $\{(32,16),(64,32)\}$ and Adam step sizes are $\{10^{-3}, 3\cdot 10^{-4}\}$.
The lag is chosen from $\{6,12\}$ (monthly), $\{8,12\}$ (weekly), or $\{14,28\}$ (daily grains).
Training minimizes squared error with early stopping on a validation slice.
The configuration with the lowest validation sMAPE, after the inverse scaling, is refit on train and validation.
Test forecasts are one-step, using observed lags.

\paragraph{M6: adaptive random forest.}
The adaptive random forest grows an ensemble of trees online and replaces trees whose error drifts~\citep{gomes2017arf}.
The cited method is a classifier for evolving streams.
We use the regression form of the same adaptive-forest mechanism, as implemented for data streams, and we do not claim the 2017 experiments as revenue results.
The forest is warmed up on history and then alternates prediction and learning at each step.
The number of trees and the feature subsample are chosen by validation sMAPE.

\paragraph{M7: Gaussian process.}
A Gaussian process prior $f \sim \mathcal{GP}(0, k)$ on the tabular covariates supplies a predictive mean and a predictive variance~\citep{rasmussen2006gp}.
On the monthly panel the kernel is chosen among radial basis, Mat\'ern, and rational quadratic, together with a noise level $\alpha \in \{10^{-2}, 10^{-1}\}$, by walk-forward validation sMAPE.
Daily fits use a radial-basis kernel and a training subsample of at most 200 points, because exact inference scales cubically in the sample size.
A fit that returns no finite forecast is omitted from the rank.

\paragraph{M8: kernel ridge.}
Kernel ridge solves~\eqref{eq:ridge} in the reproducing kernel Hilbert space of a radial-basis or Laplacian kernel, using the same tabular covariates as M1~\citep{murphy2012ml}.
Those covariates are standardized on the training block.
The kernel, the ridge penalty, and the kernel width are chosen by validation sMAPE.
Fits with test sMAPE of at least 180 are excluded from the ranking as degenerate.

\paragraph{M9: TimesFM 3.0, zero-shot.}
TimesFM is a pretrained decoder-only forecaster~\citep{das2024timesfm}.
We use the TimesFM~3.0 checkpoint \texttt{google/timesfm-3.0-pytorch} with no fine-tuning~\citep{jain2026timesfm3}.
The context window is the concatenation of the train and validation blocks.
The horizon is the test length.
The point forecast is the model's predictive mean.
If the checkpoint cannot be loaded, the family is omitted for that grain rather than replaced by another model.

\subsection{Stage 2: selective energy and generative forecasts}
\label{sec:energy}

Stage~2 asks whether a selective rule, or a conditional generative sample, improves on the Stage~1 point forecasts.
Three of the four rows are easy to mislabel.
Only E1 is TEM.

\paragraph{E1: selective TEM.}
An energy-based model defines a scalar $E_{\theta}(x, y)$ on a context $x$ and a candidate value $y$, with
\begin{equation}
  p_{\theta}(y \mid x) \propto \exp\bigl(-E_{\theta}(x, y)\bigr).
  \label{eq:ebm}
\end{equation}
Low energy is high plausibility under this model~\citep{brusokas2025tem}.
TEM trains a deterministic encoder--decoder together with such an energy head and, at inference, may reject a forecast.
Our primary backbone is TimesNet~\citep{wu2023timesnet}, trained with the published TEM routine (\texttt{Exp\_Main\_Energy}): a supervised forecast step followed by an energy step.
The series is presented as a univariate target with the same train, validation, and test borders as Table~\ref{tab:protocol}.
The input length is 12 months (label length 6) on the monthly panel and 28 days (label length 14) on the daily grains, with a one-step prediction length.
The monthly fit uses on the order of 15 supervised epochs and 10 energy epochs.
Daily fits are capped near 3 epochs of each kind.
If the TimesNet energy fit does not complete, we fall back to the published MLP predictor--decoder energy model.
That fallback is still TEM.
The candidate set at each time is the seasonal naive, the histogram-gradient-boosting forecast, and the TimesNet forecast when that forecast is available.
Let $\mathcal{C}_t$ be that set and let $\tau$ be a threshold chosen on the validation block to control validation sMAPE.
The selective forecast is
\begin{equation}
  \hat y_t =
  \begin{cases}
    \arg\min_{c \in \mathcal{C}_t} E_{\theta}(x_t, \hat y_t^{(c)})
      & \text{if }\min_{c} E_{\theta}(x_t, \hat y_t^{(c)}) \lt \tau, \\[4pt]
    y_{t-k}
      & \text{otherwise,}
  \end{cases}
  \label{eq:tem}
\end{equation}
where $k$ is the seasonal lag of the grain.
The second branch is abstention: the instrument returns the seasonal naive instead of a low-confidence energy choice.
Coverage, the fraction of test times that take the first branch, belongs with any operational reading of E1.
The sMAPE tables below do not by themselves state that coverage.

\paragraph{E2: disagreement proxy.}
E2 uses the same accept-or-abstain shape and no energy model.
Define the disagreement
\begin{equation}
  E_t = \bigl|\hat y_t^{\mathrm{HGB}} - \hat y_t^{\mathrm{seas}}\bigr|.
  \label{eq:proxy}
\end{equation}
This quantity is large when boosting and the seasonal naive disagree.
It is not the energy~\eqref{eq:ebm}.
On the validation block we set $\tau$ to an empirical quantile of $E_t$ and search the quantile levels $\{0.5, 0.6, 0.7, 0.8, 0.9\}$ by validation sMAPE.
The test rule is
\begin{equation}
  \hat y_t =
  \begin{cases}
    \hat y_t^{\mathrm{seas}} & \text{if } E_t \gt \tau, \\
    \hat y_t^{\mathrm{HGB}} & \text{otherwise.}
  \end{cases}
  \label{eq:proxy-rule}
\end{equation}
High disagreement abstains to the seasonal naive.
Low disagreement keeps the boosting forecast.
The row checks whether the TEM code is importable and then ignores it.
E2 is an ablation of the selection rule.
We do not call it TEM, and we do not identify it with E1.

\paragraph{E3: residual diffusion.}
TimeGrad forecasts a multivariate series by an autoregressive denoising diffusion process conditioned on a recurrent state~\citep{rasul2021timegrad}.
E3 keeps the denoising-diffusion step and drops the recurrent conditioner.
A three-layer network $\epsilon_{\theta}$ reads a standardized lag window, a noisy scalar target, and a normalized diffusion time, and it predicts the noise.
The forward corruption on the standardized target $y$ at diffusion index $k$ is
\begin{equation}
  y^{(k)} = \sqrt{\bar\alpha_k}\, y + \sqrt{1 - \bar\alpha_k}\,\epsilon,
  \qquad \epsilon \sim \mathcal{N}(0, 1),
  \label{eq:ddpm}
\end{equation}
and the training loss is the mean squared error between $\epsilon$ and $\epsilon_{\theta}$.
The point forecast is the mean of a small number of DDPM reverse samples, mapped back to the original scale.
Monthly fits use lag~12, 20 diffusion steps, and 8 samples.
Daily fits use lag~28, 15 diffusion steps, 4 samples, and a subsample of the training windows.
Weekly fits use lag~12, 15 steps, and 4 samples.
This is a reduced univariate sampler in the sense of TimeGrad, not the published multivariate implementation.

\paragraph{E4: score matching and Langevin sampling.}
Score-based generative models estimate $\nabla_y \log p(y)$ by denoising score matching and sample with Langevin dynamics~\citep{song2019score}.
E4 estimates a conditional score of the next standardized observation given its lag window, at one noise scale $\sigma = 0.5$.
The regression target at a corrupted point $y + \sigma\epsilon$ is $-\epsilon / \sigma^2$.
Sampling starts from Gaussian noise and repeats the Langevin update
\begin{equation}
  y \leftarrow y + \eta\, s_{\theta}(x, y) + \sqrt{2\eta}\,\xi,
  \qquad \xi \sim \mathcal{N}(0, 1),
  \label{eq:langevin}
\end{equation}
with step $\eta = 0.05$.
The point forecast is the mean of a few independent chains, inverse-scaled.
Monthly fits use 30 steps and 6 chains.
Daily and weekly fits use 20 steps and 4 chains, with the same lags as E3.
As with E3, the network is small and the fit is local to this series.

\subsection{Reproducibility}
\label{sec:repro}

This paper states the scientific protocol: grains, splits, metrics, model classes, and the selection and inference rules above.
The implementation that produced the freeze, including the full grids, epoch caps, and per-grain fit logs, is the companion notebook
\path{projects/monitoramento/notebooks/articles/arrecadacao_forecast_paper_v1.ipynb}.
Ranked metrics and stored predictions for \texttt{extracao=2026-03} are under
\path{projects/monitoramento/output/articles/arrecadacao_forecast_paper_v1/}.
The figures in Section~\ref{sec:discussion} are written to
\path{projects/monitoramento/output/figures/articles/arrecadacao_forecast_paper_v1/}
and copied into \path{figures/} beside this manuscript.
TEM training calls the published Time-Energy Model code (\url{https://github.com/JonasBrusokas/Time-Energy-Model}).
TimesFM~3.0 is the public checkpoint \texttt{google/timesfm-3.0-pytorch}.
Numbers in Section~\ref{sec:results} are read from those frozen tables.
They are not recomputed in the manuscript source.
